\subsection{Eighth Letter}

\begin{quotex}
This letter was written by \textbf{Rene Guenon} to \textbf{Guido de Giorgio} shortly after the death of
Guenon's wife. After some personal reflections, Guenon returns to some of the common themes in
earlier letters: the Sufis in Tunisia, the \emph{Ashtavakra Gita}, and, of course, some gossip about \textbf{Julius
Evola}. We see the mutual interest of the two men in \textbf{Dante}. We also see that Guenon and de Giorgio were in
complete agreements on esoteric topics, even without consulting with each other before hand. This, I hope, demonstrates
to everyone the importance of understanding sound principles; once so understood, the rest just follows. As Guenon
suggests, ``it is not the result of chance." Obviously, he would have no sympathy for the ``Infinite Monkey Theorem."

The book that Guenon referred to is most likely \emph{Crisis in the Modern World}, which had just been published in
1927. 

\end{quotex}
Blois, 6 April 1928

I am embarrassed to be so late with you; I had wanted to thank you again sooner for the expression of sympathy that you
gave me in these sad circumstances, but I didn't have the energy to write to anyone for a long
time. Your last letter arrived in Paris, and I had wanted to respond right away, but I couldn't
find even a few minutes, all my time was taken by lessons and other less interesting, although necessary, tasks.
I'm benefitting from the few days that we spent here to finally write to you; we were obliged to
come to Blois on business, otherwise, we would probably not have decided to leave Paris at this time.

I was going to thank you for your parcel that I received with great pleasure when the misfortune happened, and that in
the most sudden and least expected way: a few days of illness at most, but an illness against which they can do
nothing, cerebral-spinal meningitis. There was then a veritable epidemic in Paris: some people were taken in a few
hours! Since then, my health hasn't been very good, and there is nothing astonishing about that
after such a hard blow. However, I needed to take up again almost right away some forced tasks, but there is a little
else that I was able to do up until now.

We will return to Paris at the end of next week, that is, the 13th or 14th of April; perhaps it would be better for you
to wait until then to send me Taillard's letter which you told me about and which will surely
interest me. On my return to Paris, I will have to look again for the summary of the doctrine of the Alawis that he
translated, in order to send you a copy of it as I had promised you.

All the same, one of these days I will have to write to Evola whose several letters I have always left without response.
As of now, I have not received his \emph{Pagan Imperialism}, although he told me that he would send it to me some time
ago. Perhaps he will do it after I write him. In any case, if I receive nothing, I will tell you, and then you will be
able to lend me your copy as you suggested.

I was able to read the \emph{Ashtavakra Gita} a little quickly; it seems quite good to me. Even in this adaptation which
is surely imperfect (all the more so that the preface, as I remember, does not show much comprehension by its author);
I will have to resume it again more attentively when I have a little free time.

I received the first number of the second year of \emph{Ur}; has another one appeared since? I did not notice, on the
whole, a very big change from the first year, except that those histories of the ``magical chain" seem to take on more
and more importance, and also the ``Science of the I" disappeared from the title. Is it the ``psychological"
interpretation which I had told you about that influenced Evola to suppress that expression, which he nevertheless
seemed to be attached to? I considered all the reflections that you made in one of your letters from the beginning of
January to be spot on, regarding the subject of the preceding issue. When I write to Evola, you can be sure that I
would take account of what you said to me then, for I am convinced that you are totally correct.

Were you able to see the new work on Dante that you had spoken to me about? If yes, you will be quite kind, some time,
to tell me a little about what it is and of what type of symbolism is in it.

I carried out, some time ago, the favour which you asked me to do for Bossard; I was going to forget to tell you and to
send you the note, postmarked on the payment date.

I was happy to see that you are completely in agreement with everything that I wrote in my last book\footnote{Most likely \emph{Crisis in the Modern World}, published in 1927.}; furthermore, the contrary would have really astonished me. The
coincidences that you pointed out to me with the things that you yourself had written earlier are actually very
remarkable; these agreements are certainly not the result of chance (besides which, I don't
believe at all).

I hope that I will have one day the time to take up all your recent letters and respond at least to the most important
of the diverse questions that you pose.

I have at this moment a very disagreeable task: it is the correction of the English translation of \emph{Man and his
Becoming}; that translation was horribly done, there are some mistranslations on every page. Furthermore, instead of
sending me the manuscript as I had counted on, they typeset it and only sent me the proofs, with the result that the
editors were very unhappy when I wrote them that there were many changes to make; that is nevertheless not my fault if
they did something silly.

I like to believe that your health is not too bad now and that you haven't had another attack since
your last letter.



\flrightit{Posted on 2012-10-08 by Cologero }

\begin{center}* * *\end{center}

\begin{footnotesize}\begin{sffamily}



\texttt{Cologero on 2012-10-09 at 12:28 said: }

An accessible rendition of the Ashtavakra Gita is ``A Duet of One" by Ramesh S. Balsekar, a follower of Nisargadatta. His
commentary is antinomian, even somewhat Spinozist.


\end{sffamily}\end{footnotesize}
